%!TEX TS-program = xelatex
%!TEX encoding = UTF-8 Unicode
% Awesome CV LaTeX Template for CV/Resume
%
% This template has been downloaded from:
% https://github.com/posquit0/Awesome-CV
%
% Author:
% Claud D. Park <posquit0.bj@gmail.com>
% http://www.posquit0.com
%
% Template license:
% CC BY-SA 4.0 (https://creativecommons.org/licenses/by-sa/4.0/)
%

%-------------------------------------------------------------------------------
% CONFIGURATIONS
%-------------------------------------------------------------------------------
% A4 paper size by default, use 'letterpaper' for US letter
\documentclass[11pt, a4paper]{awesome-cv}

% Configure page margins with geometry
\geometry{left=1.4cm, top=.8cm, right=1.4cm, bottom=1.8cm, footskip=.5cm}

% Specify the location of the included fonts
\fontdir[fonts/]

% Color for highlights
% Awesome Colors: awesome-emerald, awesome-skyblue, awesome-red, awesome-pink, awesome-orange
%                 awesome-nephritis, awesome-concrete, awesome-darknight
%\colorlet{awesome}{awesome-skyblue}
% Uncomment if you would like to specify your own color
\definecolor{awesome}{HTML}{0466c8}

% Colors for text
% Uncomment if you would like to specify your own color
% \definecolor{darktext}{HTML}{414141}
% \definecolor{text}{HTML}{333333}
% \definecolor{graytext}{HTML}{5D5D5D}
% \definecolor{lighttext}{HTML}{999999}

% Set false if you don't want to highlight section with awesome color
\setbool{acvSectionColorHighlight}{true}

% If you would like to change the social information separator from a pipe (|) to something else
\renewcommand{\acvHeaderSocialSep}{\quad\textbar\quad}

%-------------------------------------------------------------------------------
%	PERSONAL INFORMATION
%	Comment any of the lines below if they are not required
%-------------------------------------------------------------------------------
% Available options: circle|rectangle,edge/noedge,left/right
%\photo[circle,noedge,right]{./profile}
\name{Cheick Tidiane}{Ba}
%\position{Resume for PHD Internship @ Twitter}%{\enskip\cdotp\enskip} 
%\position{CV}
%\address{Milan, Italy}

%\mobile{(+39)3930600691}
\email{cheick.ba@outlook.com}
\homepage{back7.github.io}
\linkedin{cheickba7}
%\twitter{@CheickBa7}
% \skype{skype-id}
% \reddit{reddit-id}
% \medium{madium-id}
\googlescholar{xRpvJ78AAAAJ&hl}{Cheick Ba}
%% \firstname and \lastname will be used
% \googlescholar{googlescholar-id}{}
\extrainfo{Orcid: 0000-0002-4035-7464 ~~~|~~~ Scopus Author ID: 57210572544 ~~~|~~~ Nazionalità: Italiana ~~~~|~~~~ Lingue: Italiano, Inglese, Francese ~~~|~~~ }

%\quote{``Be the change that you want to see in the world."}

%-------------------------------------------------------------------------------
%	BIBLIOGRAPHY
%-------------------------------------------------------------------------------
\addbibresource{references.bib}

\usepackage{multicol}

%-------------------------------------------------------------------------------
\begin{document}

% Print the header with above personal informations
% Give optional argument to change alignment(C: center, L: left, R: right)
\makecvheader[C]

% Print the footer with 3 arguments(<left>, <center>, <right>)
% Leave any of these blank if they are not needed
\makecvfooter
  {\today}
  {Cheick Tidiane Ba~~~·~~~CV~~~·~~~Made in LaTeX}
  {\thepage}

%\cvsection{Summary}
% 3rd year Ph.D. student in Computer Science at the University of Milan. Research and programming activities have led to several projects and publications, mainly in the fields of blockchain analysis, temporal and multilayer networks, and machine learning on graphs. My experience as a teaching assistant, private tutor, and volunteer developed leadership, teaching, and training skills.
%My research is in data mining and knowledge discovery, network science, and applied machine learning and statistics; with a focus on network evolution and prediction tasks on large-scale temporal data through effective and scalable computational methods.

\begin{refsection}

%-------------------------------------------------------------------------------
%	CONTENT
%-------------------------------------------------------------------------------
\cvsection{Esperienza}
\begin{itemize}
    \item \textit{Nov. 2024 – Presente}: \textbf{Pometry Ltd.}, Londra. \textit{Data Scientist}  
    Graph analytics, machine learning, data engineering e visualizzazione su larga scala — applicati a blockchain, reti finanziarie e dati complessi su grafi.

    \item \textit{Nov. 2023 – Nov. 2024}: \textbf{Queen Mary University of London}. \textit{Postdoctoral Research Assistant}
    Analisi cross‑chain utilizzando reti temporali multidimensionali con il \underline{\href{https://networks.eecs.qmul.ac.uk/}{Networks Research Group}}.

    \item \textit{Nov. 2020 – Gen. 2024}: \textbf{Università degli Studi di Milano}, Milano. \textit{Dottorando di Ricerca.}
    Reti temporali multilayer, analisi blockchain e machine learning su grafi. Responsabile della linea di ricerca su reti temporali e blockchain presso il \underline{\href{https://connets.di.unimi.it/}{CONNETS Lab}}.

    \item \textit{Mag. 2020 – Nov. 2020}: \textbf{Università degli Studi di Milano}. \textit{Ricercatore}  
    Progetto \href{https://www.vasariartexperience.it/}{VASARI}; segmentazione clienti e analisi del rischio per un gruppo finanziario italiano con il \underline{\href{https://connets.di.unimi.it/}{CONNETS Lab}}.  
\end{itemize}



\cvsection{Competenze}

%-------------------------------------------------------------------------------
%	CONTENT
%-------------------------------------------------------------------------------
\begin{cvskills}
%---------------------------------------------------------
  \cvskill
    {Programmazione} % Category
    {Python (Avanzato), Rust (Intermedio), SQL, Java, PHP, Perl, Latex} % Skills
    

%---------------------------------------------------------
  \cvskill
    {Database} % Category
    {Spark, MongoDB, Postgres, Oracle, MySQL} % Skills   
%---------------------------------------------------------
  \cvskill
    {Machine learning} % Category
    {Scikit-learn, PyTorch, Deep graph library, PyTorch geometric} % Skills  

%---------------------------------------------------------
  \cvskill
    {Data mining e modellazione dei grafi} % Category
    {Raphtory, pandas, statsmodels, networkx, networkit, Gephi} % Skills  

    %---------------------------------------------------------
  \cvskill
    {Sviluppo Web} % Category
    {MongoDB, Express, REST API, React, AngularJS, HTML5/CSS/JavaScript} % Skills

 % %---------------------------------------------------------
 %  \cvskill
 %    {DevOps} % Category
 %    {Git, Jupytherhub} % Skills

%---------------------------------------------------------
  \cvskill
    {Lingue} % Category
    {Inglese (C2), Italiano (madrelingua), Francese (madrelingua), Spagnolo (base), Tedesco (base)} % Skills
    
    \cvskill
    {Extracurricular} % Category
    {Basket, escursioni, nuoto, viaggi, apprendimento delle lingue, lettura}
\end{cvskills}

\cvsection{Educazione}

\begin{itemize}
\item \textit{26 Gen. 2024}: \textbf{Università degli Studi di Milano, Italia} – \textit{PhD in Informatica} - cum laude, \textit{Doctor Europaeus}.


    \item \textit{11 Mar. 2020}: \textbf{Università degli Studi di Milano, Italia} - \textit{Laurea Magistrale in Informatica}, 110/110 cum laude.
\end{itemize}

\begin{itemize}
    \item \textit{Mag. 2023 – Ago. 2023}: \textbf{Queen Mary University of London} – \textit{PhD Student in visita}

    \item \textit{Mag. 2022 – Ago. 2022}: \textbf{CIRAD Montpellier} - \textit{PhD Student in visita}
    % Topics: machine learning on multilayer networks, text mining for food security
\end{itemize}

%----------------------------------------------------------------------------------------------------------------------------------



%----------------------------------------------------------------------------------------------------------------------------------


%-------------------------------------------------------------------------------
%	SECTION TITLE
%-------------------------------------------------------------------------------


\cvsection{Pubblicazioni}    

Ho contribuito a 23 pubblicazioni peer-reviewed, tra cui riviste e conferenze internazionali. Tra queste, 8 sono in riviste Q1. Un elenco completo dei miei contributi è disponibile su \underline{\href{https://scholar.google.com/citations?user=xRpvJ78AAAAJ&hl=it}{Google Scholar}}, \underline{\href{https://dblp.uni-trier.de/pid/247/1543.html}{DBLP}}, e \underline{\href{https://www.scopus.com/authid/detail.uri?authorId=57210572544}{Scopus}}.  

\cvsection{Presentazioni}

Ho presentato la mia ricerca in oltre 20 conferenze principali sulla network science, sistemi complessi e machine learning, tra cui il Mathematical Institute di Oxford, NetSci, CCS, ECMLPKDD, WebSci. I miei interventi si sono concentrati su argomenti come l'analisi di grafi temporali multidimensionali, data mining, blockchain e sistemi decentralizzati.


\cvsection{Premi \& Riconoscimenti}

\begin{itemize}
    \item \textbf{2025}: \textbf{FOSS Award}, 20° International Workshop on Mining and Learning with Graphs – ECMLPKDD.
    \item \textbf{2023}: \textbf{Premio “miglior poster”}, 20° International Workshop on Mining and Learning with Graphs – ECMLPKDD.
    \item \textbf{2021}: \textbf{Premio “Runner‑up for Best Paper”}, 2° posto su 140 sottomissioni alla conferenza GoodIT ‘21.
\end{itemize}

\cvsection{Insegnamento}


\begin{itemize}
    \item \textbf{2022/2023}, \textbf{Università di Pavia, Italia} – \textit{Assistente Didattico}, Intelligenza Artificiale (12 ECTS), 24h  
    \item \textbf{2021/2022, 2022/2023}, \textbf{Università degli Studi di Milano, Italia} – \textit{Docente}, Social Media Mining (12 ECTS), 7h  
    \item \textbf{2020/2021, 2021/2022}, \textbf{Università degli Studi di Milano, Italia} – \textit{Assistente Didattico}, Analisi delle Reti Sociali (6 ECTS), 20h  
\end{itemize}


\end{refsection}
\end{document}
